\section{The full polar-retaining rational matrix under a right zero-free half-plane}
\label{sec:polar-common-matrix}

The positive source-column result MA-MO-016 does not by itself
extend the original matrix on the complementary diagonal source.
Here that second column is constructed under an explicit
zero-location premise. Both original inverses are evaluated on
the positive narrow line, and both polar moments are retained.
In particular a zero on the critical line does not obstruct this
construction when the stipulated premise A is available.

Fix
\[
 \begin{gathered}
 0\le a_*<a<1/2,\qquad s\ge0,\\
 \zeta(\sigma+it)\ne0
          \quad\hbox{whenever }\sigma>1/2+a_*.
 \end{gathered}
 \tag{MA-MO.320}
\]
The inequality refers to zeros, not to the pole at one. It is
an additional arithmetic hypothesis in this section; A gives
it with \(a_*=0\). Keep the rational place set
\(S=\{\infty\}\), the actual cutoff \(\phi\), and the
Gaussian sections of MA-MO-016. The cutoff need not be symmetric.
Put
\[
 \begin{aligned}
 H&=H_{\Sigma,a,s},&H_a&=H_{a,s},&
 \mathcal W&=H\oplus\mathbb C^2,\\
 q(x)&=1-\phi(x),&p_-(x)&=e^{-x/2},&
 D&=P_+-P_- .
 \end{aligned}
 \tag{MA-MO.321}
\]
In this section \(p_-(x)\) denotes the displayed scalar
polar function; the source retraction remains denoted
\(p_-^{\rm src}\) when it is needed. The operator \(P_-\)
is always the original source projection.

\paragraph{MA-MO-018 (actual full common-space trace-class extension).}
For every \(g\in C_c^\infty(\mathbb R)\), the original
smooth operator
\[
       C_gD:\mathscr S_{\rm sum}\longrightarrow
                                     \mathscr S_{\rm sum}
\]
extends uniquely to a trace-class endomorphism
\[
       \mathscr D_g:\mathcal W\longrightarrow\mathcal W.
 \tag{MA-MO.322}
\]
Its restriction to the positive completion is the actual
source-column map of MA-MO-016. On each fixed compact test
support, \(g\mapsto\mathscr D_g\) is continuous into trace
class. This is an extension of the original wide-inverse
source action, with no substitution of an orthogonal projector
or of an unrelated continued inverse.

The trace and Fredholm determinant of (MA-MO.322) are therefore
defined. The theorem does not compute that trace, assert that
\(\mathcal W_+\) is invariant under \(\mathscr D_g\), or
define its action on \(\mathcal W/\mathcal W_+\). The exact
outgoing condition of (MA-MO.256) remains in force. In
particular this is not a zero-location-free proof of a full
narrow identity.

\begin{proof}
\emph{The original inverse on every negative smooth input.}
For \(y\in\mathcal E\), let
\[
 F_y(x)=\sum_{n\ge1}\mu(n)n^{-1/2}y(x+\log n).
 \tag{MA-MO.323}
\]
This is the original wide Euler inverse. The series and its
derivatives converge locally absolutely, by the rapid right
decay of \(y\). On a line \(b>1/2\), it converges in the
weighted Hilbert operator norm and has Mellin transform
\(\mathcal My(z)/\zeta(1/2+z)\), as in MA-MEY-001.

Hypothesis (MA-MO.320) gives a fixed horizontal margin
between every line \(\Re z\ge a\) and every nontrivial
zero. The reciprocal minimum-modulus proof of MA-RS-004
therefore gives polynomial bounds for the reciprocal and
all its derivatives, uniformly on each fixed closed vertical
strip in \(\Re z\ge a\). The inverse is holomorphic there:
the zeta pole at \(z=1/2\) gives a zero of the inverse,
not a pole. Thus contour shifting in the rapidly decreasing
Mellin quotient is legitimate. It proves
\[
 \begin{gathered}
 e^{bx}F_y\in\mathscr S(\mathbb R)\quad\hbox{for every }b\ge a,
 \qquad
 V_aF_y=j_aV_ay,\\
 A(z)=\zeta(1/2+z),\qquad j_a(t)=A(a+it)^{-1},\qquad
 \mathcal MF_y(1/2)=0.
 \end{gathered}
 \tag{MA-MO.324}
\]
Here \(V_a=\mathcal F_+M_{e^{ax}}\) is the unitary
Fourier normalization. The last Mellin integral converges
absolutely: use a line \(a<1/2\) at the left and a line
\(b>1/2\) at the right. Its value is zero because the
holomorphic quotient vanishes at \(z=1/2\). The equality
with the original function (MA-MO.323) follows by starting
on the absolutely convergent wide line before shifting.
This argument is for every \(y\in\mathcal E\), with no
ideal membership assumption.

Write \(J_a\) for the maximal inverse multiplier on
\(H_{a,0}\). Then (MA-MO.324) says precisely
\(F_y=J_ay\) on this actual source core. The same is true
for \(Ry\). No contour is moved to the negative narrow
line in making either inverse identification.

\emph{The two actual source functions and their polar moments.}
Define an even additive function by
\[
 u_y(t)=\frac{q(\log|t|)F_y(\log|t|)}{2\sqrt{|t|}}
        \quad(t\ne0),\qquad u_y(0)=0.
 \tag{MA-MO.325}
\]
It is Schwartz: it vanishes near additive zero, and all
its derivatives decrease rapidly at infinity by the
arbitrarily large right weights in (MA-MO.324).
The exact source formula for \(p_+\) on a diagonal vector
is consequently
\[
        p_+i_-y=h_y:=u_y+\mathfrak F u_{Ry},
             \qquad h_{Ry}=\mathfrak F h_y.
 \tag{MA-MO.326}
\]
This is the original map of Meyer's
\texttt{lem:Hilm\_p\_embed}, specialized to the rational
even sector. It uses the wide inverse (MA-MO.323), now
identified with its positive narrow multiplier value.

Let \(\lambda_0,\lambda_1\) be the bounded functionals
of MA-MO-011 on \(H\). For arbitrary \(y\in\mathcal E\),
not only ideal vectors, the same bilinear Fourier computation
and (MA-MO.324) give
\[
 \lambda_0(y)=\int e^{x/2}\phi(x)F_y(x)\,dx,
                  \qquad\lambda_1(y)=\lambda_0(Ry).
 \tag{MA-MO.327}
\]
The functionals are bounded on \(H\), since the local
weighted step is Schwartz and \(j_a\) has polynomial
growth on this line. The zero total moment in (MA-MO.324)
now proves both source moment identities
\[
       h_y(0)=-\lambda_1(y),\qquad
       \int h_y=-\lambda_0(y).
 \tag{MA-MO.328}
\]
Indeed \(\int u_y=\int e^{x/2}qF_y=-\lambda_0(y)\),
while \(u_y(0)=0\); Fourier evaluation and integration
at zero give the other term. Thus neither source polar
coefficient has been discarded.

Let \(G(z)=\mathcal Mg(z)\). Under the actual additive
dilation operator \(\mathcal D_g\) from (MA-MO.263),
the two moments are multiplied by \(G(-1/2)\) and
\(G(1/2)\), respectively. Hence the polar coordinates
of \(C_gDi_-y\) are
\[
 \nu_g(y)=
       \bigl(-G(-1/2)\lambda_1(y),
             -G(1/2)\lambda_0(y)\bigr).
 \tag{MA-MO.329}
\]
This is already a bounded finite-rank map \(H\to\mathbb C^2\).

\emph{An explicit positive-line off-diagonal kernel.}
Let \(\mathcal H\) be the radial form of the actual even
additive Fourier transform, so that
\[
 \mathcal M\mathcal H F(z)=m(z)\mathcal MF(-z),\qquad
 m(z)=\frac{\Gamma_{\mathbb R}(1/2+z)}
                    {\Gamma_{\mathbb R}(1/2-z)}.
 \tag{MA-MO.330}
\]
Only \(z=a+it\) is used in this step. Put
\[
       P(\xi)=\int e^{\xi x}\phi'(x)\,dx,
                  \qquad z=a+it,\quad w=a+iu.
\]
For \(\Re\xi<0\), integration by parts, with zero
boundary terms because \(q=0\) far left, gives
\(\int e^{\xi x}q(x)\,dx=P(\xi)/\xi\).
It follows directly by Fourier inversion that the source
operator \(C_gZ_a\mathcal H M_qJ_a\) has kernel
\[
       o_{a,g}(t,u)=\frac{G(z)A(z)m(z)j_a(u)}{2\pi}
                          \frac{P(-z-w)}{-z-w}
 \tag{MA-MO.331}
\]
in the \(V_a\) coordinates. Here \(Z_a\) is the maximal
multiplier \(A(a+it)\). On the source core its sequential
domain is valid: \(qF_y\) comes from an additive Schwartz
function, its local Fourier transform has weighted
Schwartz estimates at \(a>-1/2\), and the final zeta
multiplier has polynomial growth. The preceding inversion
can be justified absolutely after inserting the Schwartz
input transform; the separated real part
\(\Re(-z-w)=-2a\) is crucial.

All derivatives of (MA-MO.331) are Schwartz in \((t,u)\).
The factors \(A,m,j_a\) and their derivatives have
polynomial growth on their fixed lines. The two rapid
factors are \(G(a+it)\) and \(P(-2a-i(t+u))\), and
\(|u|\le |t|+|t+u|\). The denominator stays a distance
\(2a\) from zero. Thus (MA-MO.331) extends to a
trace-class operator
\(O_a(g):H_{a,s}\to H_{a,s}\) for every \(s\ge0\),
by the Bessel/Hermite argument of MA-MO.126. This is an
extension of the actual sequential operator on the specified
smooth core, not a claim that all four unbounded factors
are separately defined on the whole Hilbert space.

\emph{The remaining half-line multiplier uses the sum norm.}
Let \(K_a(g)\) be the trace-class operator of MA-MO.125,
which agrees on \(\mathcal E\) with
\(C_g(M_\phi-Z_aM_\phi J_a)\). The agreement here holds
on all of \(\mathcal E\): (MA-MO.324) makes the inverse
transform Schwartz, and compact frequency cutoff with
\(j_a\chi_NV_ay=\chi_NV_aF_y\) gives the same graph
approximation as in MA-MO.128--129.

The additional map
\(N_a(g)=C_gM_\phi:H\to H_{a,s}\) is trace class as
well. After the two weight isometries, its physical kernel is
\[
 \frac{e^{ax}(1+x^2)^{s/4}}
      {(e^{2ay}+e^{-2ay})^{1/2}(1+y^2)^{s/4}}
                         g(x-y)\phi(y).
 \tag{MA-MO.332}
\]
The difference \(x-y\) lies in a fixed compact set. At
\(y\to-\infty\) the exponential ratio decreases like
\(e^{2ay}\), up to a bounded factor depending on that
compact set; at the other end \(\phi(y)\) vanishes.
The same estimates hold after every derivative, and the
polynomial weight ratio is harmless on a bounded difference.
Thus (MA-MO.332) is Schwartz in \((x,y)\), and its Hermite
expansion proves the assertion. No trace-class claim for
\(C_gM_\phi\) on a single unchanged one-sided space is
being made.

\emph{Removing the actual polar sections before combining the rows.}
Write \(f_0,f_1\) for (MA-MO.251), and set
\[
       e_0=\Theta f_0,\qquad
       e_1=\Theta f_1-p_-,\qquad Re_0=e_1.
 \tag{MA-MO.333}
\]
Both vectors belong to \(H\). The theta estimate at the
right and Poisson summation at the left show that \(e_0\)
has left tail \(-e^{x/2}\), while \(e_1\) has right tail
\(-e^{-x/2}\), with rapid opposite tails. These are
integrable in the sum weight since \(a<1/2\). The reflection
identity follows from \(\mathfrak Ff_0=f_1\) and (MA-MO.84).

For any even Schwartz source \(h\), the exact centered
pole subtraction gives, in locally integrable functions,
\[
       \Theta h=Z_aF_h+\left(\int h\right)p_-,
                    \qquad F_h(x)=2e^{x/2}h(e^x).
 \tag{MA-MO.334}
\]
This is the contour identity of MA-MO-004: shifting the
wide direct theta transform to \(a<1/2\) crosses only
the zeta pole, with the displayed positive residue. Zeta
zeros are not poles in this direct summation step.

Subtract the two actual sections from \(C_gDi_-y\)
using its moments (MA-MO.329). The first component of the
remaining diagonal vector is exactly
\[
 \begin{aligned}
 T_{a,g}y={}&K_a(g)y-N_a(g)y+O_a(g)Ry\\
       &+G(-1/2)\lambda_1(y)e_0
         +G(1/2)\lambda_0(y)e_1.
 \end{aligned}
 \tag{MA-MO.335}
\]
To check the formula, the radial vector of \(h_y\) is
\(qF_y+\mathcal H(qF_{Ry})\). Apply (MA-MO.334) and
(MA-MO.328), and use
\(Z_aM_qJ_a-1=(M_\phi-Z_aM_\phi J_a)-M_\phi\).
The subtracted source moments turn \(\Theta f_1\) into
\(\Theta f_1-p_-\), with exactly the signs in
(MA-MO.335). Every term in that formula is now a
trace-class map \(H\to H_{a,s}\): use the bounded input
inclusion for \(K_a\), (MA-MO.332) for \(N_a\), reflection
followed by that inclusion for \(O_a\), and the two
bounded functionals for the finite columns.

The second component of the same diagonal remainder is
\(RT_{a,g^\vee}Ry\), where \(g^\vee(x)=g(-x)\).
Indeed (MA-MO.326) gives \(h_{Ry}=\mathfrak Fh_y\),
Fourier scaling gives
\(\mathfrak F\mathcal D_g=\mathcal D_{g^\vee}\mathfrak F\),
and the two Gaussian sections are interchanged by Fourier.
The two components therefore agree on every
\(y\in\mathcal E\), by the original Poisson identity.
Each defines a trace-class map from \(H\) to its respective
weighted component. The diagonal isometry
\(\iota:H\to H_{a,s}\oplus H_{-a,s}\) has closed range.
Since \(\mathcal E\) is dense in \(H\), the paired map
has range in this diagonal on the entire completion.
Consequently there is a unique trace-class map
\[
       T_g:H\to H,\qquad
       \iota T_g y=(T_{a,g}y,RT_{a,g^\vee}Ry).
 \tag{MA-MO.336}
\]
This uses the actual sum norm and source compatibility;
it does not replace that norm by either one-sided norm.

\emph{The complete original matrix.}
The negative source column has now extended as
\(y\mapsto(T_gy,\nu_g(y))\in\mathcal W\), a trace-class
map. The two polar source columns
\(C_gDs_j\) are actual fixed vectors of
\(\mathscr S_{\rm sum}\subset\mathcal W\), as proved in
MA-MO-016. Thus in the source-section coordinates the full
extension is
\[
       \mathscr D_g(y,c)=
             (T_gy,\nu_g(y))
                       +\sum_{j=0}^1c_j\,C_gDs_j.
 \tag{MA-MO.337}
\]
It is trace class and agrees with the original operator on
the dense smooth domain
\(\mathfrak s(\mathcal E\oplus\mathbb C^2)\).
This proves uniqueness and its agreement with MA-MO-016
on \(\mathcal W_+\). All kernel estimates use finitely
many seminorms of \(g\) on a fixed compact support, and
the finitely many actual source columns depend continuously
on those tests. This proves the stated trace-norm continuity.
\end{proof}

\paragraph{Exact scope of the repaired matrix.}
The two positive-line inverse identities (MA-MO.324) are
where the zero-free right half-plane is used. The argument
does not infer them from a generic zero-free line alone, nor
does it transport an original inverse across unaccounted
right-exterior zeros. The two cutoff moments and actual
polar sections are required by (MA-MO.328)--(MA-MO.335).
They are not zero for a general negative source input.

The matrix (MA-MO.322) is a defined trace-class endomorphism
of the specified enlarged common completion. The positive
source endomorphism \(\jmath_+A_{+,g}\) evaluated in
MA-MO-017 is a different operator on a different space;
its finite-column difference is (MA-MO.299). Likewise
the raw full matrix need not preserve \(\mathcal W_+\),
so its quotient trace cannot be assigned by deleting that
distinction. No new arithmetic zero-location assertion,
or unconditional narrow trace identity, follows here.

The source maps are the freshly checked Meyer (2005)
\texttt{lem:Hilm\_p\_embed}, defining \(P_+,P_-\),
\texttt{eq:Poisson\_summation} and
\texttt{lem:Poisson\_sumpm}, author-manuscript pp.~39--41,
supplied \texttt{main.tex} lines~3808--4043.
The new analytic comparison is the explicit source inverse
shift and the two-row construction above. In the inherited
cote~1 nuclear framework, scans~079--080 and~111--112
through ALPOGE\_BRIDGES BR-TR-01/02, this now supplies a
defined endomorphism in the trace ideal, rather than only
a formal finite-length or exact-sequence class. It does not
by itself evaluate that endomorphism trace.
